\paragraph{A prescribed diameter.}
The joint envelope also admits an exact diameter restriction. The
diameter may be an upper bound or a prescribed value; the domain
otherwise varies.

\begin{theorem}[The joint envelope with diameter information]
\label{jointdiam:main}
Assume \eqref{boundarygap:assumptions} and let $D>0$. For $g,P>0$ put
\begin{equation}\label{jointdiam:definition}
 \begin{gathered}
 h=\min\{g,P^2/m,PD\},\qquad
 \rho_-=h^2/P^2,\qquad \rho_+=\min\{h/m,D^2\},\\
 \mathcal J_D(g,P)=\max_{\rho_-\le\rho\le\rho_+}V(h,\rho).
 \end{gathered}
\end{equation}
Set $\mathcal J_D(g,0)=\mathcal J_D(0,P)=0$. Then
\[
 G(x)\le g,\quad \|\nabla f(x)\|\le P,\quad
 \operatorname{diam}Q\le D
 \quad\Longrightarrow\quad
 \delta(x)\le\mathcal J_D(g,P).
\]
The bound is sharp for these three upper bounds. The same supremum
holds with $\operatorname{diam}Q=D$, and is approached by smooth
functions on neighborhoods of triangles with nonempty interior in
$\mathbb R^2$.

For prescribed positive values $G(x)=g$ and $\|\nabla f(x)\|=P$,
the class is nonempty exactly when $g\le PD$. Whenever it is
nonempty, its supremum is again $\mathcal J_D(g,P)$, whether the
diameter is bounded by $D$ or fixed at $D$.
\end{theorem}

For $H_1=0$ put
$\beta=(H_0^2-2mH_0-m^2)/(4H_0)$. Then
\begin{equation}\label{jointdiam:constant}
 \mathcal J_D(g,P)=
 \begin{cases}
 h+\beta\rho_+,&\beta\ge0,\\
 h+\beta\rho_-,&\beta<0.
 \end{cases}
\end{equation}
For example, $m=1$, $H_0=4$, $H_1=0$, $g=1$, $P=2$ and
$D=3/5$ give
\[
 \beta=7/16,\qquad
 \mathcal J_D(1,2)=463/400<23/16=\mathcal J(1,2).
\]
Both observations can be exact, since $g=1<6/5=PD$.
For $H_1>0$ and $\mathcal S(h)\le0$, radius monotonicity gives
\begin{equation}\label{jointdiam:ordinary}
 \mathcal J_D(g,P)=V(h,\rho_-)
       \le h+\rho_-\mathcal S(h)\le h.
\end{equation}
Thus the ordinary-gap regime still needs only one switch equation,
or the displayed algebraic upper bound.

\begin{corollary}[The norm envelope with diameter information]
\label{jointdiam:norm}
For $P>0$ let $r_D=\min\{D,P/m\}$ and $q=\max\{P-mD,0\}$. The
sharp norm-only bound with diameter at most $D$ is
\begin{equation}\label{jointdiam:norm-formula}
 \mathcal E_D(P)=\mathcal J_D(PD,P)
       =A(P,r_D)=\mathcal E_q(P).
\end{equation}
It remains sharp with both $\|\nabla f(x)\|=P$ and
$\operatorname{diam}Q=D$. Set $\mathcal E_D(0)=0$.
If $H_1=0$ and $P\ge mD$, then
\begin{equation}\label{jointdiam:norm-constant}
 \mathcal E_D(P)=PD+\beta D^2.
\end{equation}
For fixed $D$ this is an affine function of $P$.
\end{corollary}
